\subsection{Ventajas, parámetros y limitaciones del diseño}

El principal beneficio de esta arquitectura es que desacopla el ahorro de memoria de la necesidad de reconstruir por completo el vector de estado en
cada operación. Además, introduce parámetros claros y controlables: tamaño del bloque, número de entradas en caché, política de reemplazo y tipo de
compresor sin pérdida. Esto permite adaptar la arquitectura a diferentes patrones de acceso y a distintas disponibilidades de memoria.

No obstante, la arquitectura también presenta limitaciones. Si los estados cuánticos generados por el algoritmo son poco compresibles, la reducción
de memoria puede resultar modesta. Del mismo modo, si las compuertas inducen accesos dispersos sobre muchos bloques distintos, la tasa de fallos de
caché puede crecer y aumentar la cantidad de descompresiones. Por esta razón, el análisis experimental posterior debe estudiar el compromiso entre
ratio de compresión, latencia y tamaño del conjunto de trabajo.

En conjunto, los elementos anteriores delimitan un espacio de compromiso entre ahorro de
memoria, granularidad de acceso y costo temporal de compresión y descompresión. Sobre esa
base, la siguiente sección presenta la materialización de estas decisiones en un mecanismo de
almacenamiento comprimido apoyado en Blosc.
